\documentclass[12pt,a4paper]{article}
\usepackage{minekommandoer}
\begin{document}
\begin{center}
    \setfontsize{16pt}{\textbf{Løsningsforslag til prøve i støkiometri:}}
\end{center}\vspace{1cm}
\textbf{Oppg. 1:}\nl
a) beregn den molare massen til følgende stoffer:\nl
$H_2O$, $NH_3$, $H_3PO_4$,$CaCO_3$,$KMnO_4$,$Ca(ClO_3)_2$\nl[2]
Her summerer vi atommasene til alle atomene i molekylene/ionene. Verdiene finner vi periodesystemet.\nl
$Mm(H_2O)=2\cdot Mm(H)+Mm(O)=2\cdot 1,01g/mol+16,0g/mol=18,02g/mol$\nl
$Mm(NH_3)=Mm(N)+3\cdot Mm(H)=14,01g/mol+3\cdot 1,01g/mol=17,04g/mol$\nl
$Mm(H_3PO_4)=3\cdot Mm(H)+Mm(P)+4\cdot Mm(O)=3\cdot 1,01g/mol+30,98g/mol+4\cdot 16,0g/mol=98,01g/mol$\nl
$Mm(CaCO_3)=Mm(Ca)+mM(c)+3\cdot Mm(o)=\Caa+\Ca+3\cdot \Oa=100,1g/mol$\nl
$Mn(KMnO_4)=Mm(K)+Mm(Mn)+4\cdot Mm(O)=\Ka+\Mna+4\cdot \Oa=158,0g/mol$\nl
$Mm(Ca(ClO_3)_2)=Mm(Ca)+2\cdot (Mm(Cl)+3\cdot MM(O))=\Cla+2\cdot(\Ca+3\cdot\Oa)=207g/mol$\nl
Her er det $Ca(ClO_3)_2$ som er vanskelig.  Vi har en parentes her med 2 på utsiden.  Det betyr at atomene inn parentesen skal 2 ganger. altså 2 klor og $3\cdot 2=6$ oksygenatomer. 
\newpage
b) Beregn hvor mange mol du har av stoffene i a når du har 34,5g av stoffene.\nl
Når vi skal svare på denne oppgaven kan vi bruke trekanten:\nl
\trekant\nl[2]
Vi skal beregne antall mol av stoffene i a. om vi b ruker trekanten så ser vi at:\nl
$$n=\frac{m}{Mm}$$\nl
Vi får:\nl
$n(H_2O)=\frac{34,5g}{18,02g/mol}=1,91mol$\nl
$n(NH_3)=\frac{34,5g}{17,03g/mol}=2,02g/mol$\nl
$n(H_3PO_4)=\frac{34,5g}{98,01g/mol}=0,35mol$\nl
$n(CaCO_3)=\frac{34,5g}{100,1g/mol}=0,34mol$\nl
$n(KMnO_4)=\frac{34,5g}{158,0g/mol}=0,22mol$\nl
\newpage
c) Beregn hvor mange gram du har av stoffene i a når du har 14,5mol av stoffene.\nl
Om vi igjen ser på  trekanten så er det slik at:
$$m=n\cdot Mm$$\nl
Vi finner altså antall gram av stoffet ved å multiplisere antall mol med den molare massen\nl
Vi får:
$m(H_2O)=14,5mol\cdot 18,02g/mol=261,29g$\nl
$m(NH_3)=14,5mol\cdot 17,03g/mol=246,94g$\nl
$m(H_3PO_4)=14,5mol\cdot 98,01g/mol=1421,15g$\nl
$m(CaCO_3)=14,5mol\cdot 100,1g/mol=1451,45g$\nl
$m(KMnO_4)=14,5mol\cdot 158,0g/mol=2291g$\nl
$m(ca(CO_3)_2)=14,5mol\cdot 207,0g/mol=3001,5g$\nl
\nl[2]
\textbf{Oppg. 2:}\nl
Gitt følgende kjemiske reaksjon:\nl
$$C+O_2=CO_2$$\nl
a) Finn molar masse til C, $O_2$ og $CO_2$.\nl
$Mm(C)=\Ca$\nl
$Mm(O_2)=2\cdot\Oa=32,00g/mol$\nl
$Mm(CO_2)=\Ca+2\cdot \Oa=44,00g/mol$
\newpage
b) Hvor mye får vi laget av $CO_2$ og hvor mye $O_2$ trenger vi når vi har 12g C?\nl
Her må vi først se på hvor mange mol vi har når vi har 12,0g C:\nl
$n(C)=\frac{m(C)}{Mm(C)}=\frac{12g}{\Ca}=1mol \ C$\nl
Så må vi se på reaksjons-ligningen:\nl
\begin{table}[H]
    \centering
    \begin{tabular}{cccccc}
         r.x.:& C&+ &$O_2$&= & $CO_2$\\
         mol forhold:&1 &  & 1 &  &1 \\
         gram:&$1mol\cdot\Ca$ &  &$1mol\cdot 32,0g/mol$  &  &$1mol\cdot 44,0g/mol$ \\
    \end{tabular}
\end{table}\nl
Av dette oppsettet ser vi at om vi har ett mol C (12,0g) så trenger vi 1 mol av $O_2$ og ett mol $CO_2$. Det betyr at vi trenger 32,0g $O_2$ og får laget 44,0g $CO_2$. 
\newpage
c) Hvor mange gram trenger vi av C og $O_2$ for å lage 345g $CO_2$?\nl
Vi må først beregne hvor mange mol vi får laget av $CO_2$.  Vi bruker trekanten og beregner dette:\nl
$n(CO_2)=\frac{m(CO_2)}{Mm(CO_2)}=\frac{345g}{44,0g/mol}=7,84mol$\nl
Så må vi se på reaksjons-ligningen:\nl
\begin{table}[H]
    \centering
    \begin{tabular}{cccccc}
         r.x.:& C&+ &$O_2$&= & $CO_2$\\
         mol forhold:&1 &  & 1 &  &1 \\
         mol:&$1\cdot 7,84mol$ &  &$1\cdot 7,84mol$  &  &$1\cdot 7,84mol$ \\
         gram:&$7,841mol\cdot\Ca$ &  &$7,841mol\cdot 32,0g/mol$  &  &$7,841mol\cdot 44,0g/mol$ \\
         gram:&$94,092g$ & + &$250,912g=345,004$  &  &$345,004g$ \\
    \end{tabular}
\end{table}\nl
Av oppsettet over ser vi at vi trenger 94,092g C og 250,912g $O_2$ for å lage 345g $CO_2$.
\nl[2]

\noindent\textbf{Oppg. 3:}\nl
Gitt følgende reaksjon:
$$N_2+3H_2=2NH_3$$
a) hvor mye ammoniakk ($NH_3$) får vi dannet om vi starter med 32g $N_2$, og hvor mye $H_2$ trenger vi da?\nl
Vi begynner med  beregne hvor mange mol $N_2$ vi har:\nl
$n(N_2)=\frac{m(N_2)}{Mm(N_2)}=\frac{32g}{28g/mol}=1,143mol$
\newpage
Så må vi se på reaksjons-ligningen:
\begin{table}[H]
    \centering
    \begin{tabular}{cccccc}
         r.x.:& $N_2$&+ &$3H_2$&= & $2NH_3$\\
              mol forhold:&1 & &2 & &3 \\
              mol:&$1\cdot 1,143mol$ & &$3\cdot 1,143mol$ & &$2\cdot 1,143mol$ \\
              mol:&$1,143mol$ & &$3,429mol$ & &$2,286mol$ \\
              gram::&$1,143mol\cdot 28,0g/mol$ & &$3,429mol\cdot 2,02g/mol$ & &$2,286mol\cdot 17,03g/mol$ \\
    \end{tabular}
\end{table}\nl
b) Hvor mye ammoniakk får du dannet om du startrer med 32g $N_2$ og 5g $H_2$?\nl
Dette er en vanskelig oppgave. Den krever god forståelse for å få den til.  Vi starter med å beregne antall mol av $N_2$ og $H_2$.\nl
$n(N_2)=\frac{m(N_2)}{Mm(N_2)}=\frac{32g}{28g/mol}=1,143mol$\nl
$n(H_2)=\frac{m(H_2)}{Mm(H_2)}=\frac{5g}{2,02g/mol}=2,475mol$\nl
Så hva sier dette oss? Om vi ser på reaksjonen, så skal vi ha 3 gagner så mange mol $H_2$ som $N_2$.  Vi har 1,143mol $N_2$. Om  alt dette skal reagere så trenger vi $3\cdot 1,143mol=3,429mol H_2$.  Men så mange mol $H_2$ har vi ikke. Vi har 2,475mol $H_2$.  Det er det som vil reagere av $H_2$.  Deler vi dette tallet på 3 får vi antall mol $N_2$ som reagerer.\nl
Antall mol $N_2$ som reagerer:\nl
$$n(N_2)=\frac{2,475mol}{3}=0,825mol$$\nl
Antall gram $N_2$ som reagerer:\nl
$m(N_2)=0,825mol\cdot 28g/mol=23,1g$\nl
23,1g $N_2$ vil altså reagere med 5g $H_2$.  Summen av disse blir da mengden ammoniakk som dannes:  23,1g+5g=28,1g\newpage
Alternativt:\nl
Når 0,825mol $N_2$ reagerer får vi dannet $2\cdot 0,825mol=1,65mol \ NH_3$\nl
$m(NH_3)=1,65mol\cdot 17,03g/mol=28,1g$

\end{document} 